\documentclass[11pt]{article}

\usepackage[a4paper,margin=2.0cm]{geometry}
\usepackage{amsmath,amssymb,amsthm,mathtools}
\usepackage{microtype}
\usepackage[hidelinks]{hyperref}
\hypersetup{
  pdftitle={Integral Solutions of y^2+x^2y+xz^2=2},
  pdfauthor={Jamal Agbanwa}
}

\newtheorem{theorem}{Theorem}
\newtheorem{proposition}[theorem]{Proposition}
\theoremstyle{remark}
\newtheorem*{remark}{Remark}

\newcommand{\Z}{\mathbb Z}
\newcommand{\Sq}{\{r^2:r\in\mathbb Z\}}
\setlength{\parindent}{0pt}
\setlength{\parskip}{4pt}
\setlength{\abovedisplayskip}{5pt}
\setlength{\belowdisplayskip}{5pt}
\setlength{\abovedisplayshortskip}{3pt}
\setlength{\belowdisplayshortskip}{3pt}

\title{Integral Solutions of \texorpdfstring{$y^2+x^2y+xz^2=2$}{y2+x2y+xz2=2}\\[1mm]
\large Complete Classification, an Infinite Family, and the Guiding Idea}
\author{Jamal Agbanwa}
\date{}

\begin{document}
\maketitle
\vspace{-5mm}

\begin{abstract}
We give an exact elementary description of every integral solution of
$y^2+x^2y+xz^2=2$.  We then exhibit an explicit infinite family on the
single fibre $x=-17$.  The final section explains how the quadratic symmetry,
a norm equation, and a deliberately engineered Pell unit lead naturally to
the numbers occurring in that family.
\end{abstract}
\vspace{-2mm}

\section{Complete classification}

Let $\Sq$ denote the set of nonnegative integral squares.

\begin{theorem}[Exact divisor--square criterion]\label{thm:classification}
There is no solution with $x=0$.  For $x\ne0$, all integral solutions are
precisely
\begin{equation}\label{eq:classification}
 (x,y,z)=
 \left(
   \varepsilon d,\ m,\
   \delta\sqrt{-\varepsilon\left(dm+\frac{m^2-2}{d}\right)}
 \right),
\end{equation}
where
\[
 \varepsilon,\delta\in\{\pm1\},\qquad d\in\mathbb Z_{>0},\qquad m\in\mathbb Z,
\]
subject only to
\begin{equation}\label{eq:conditions}
 d\mid m^2-2,
 \qquad
 -\varepsilon\left(dm+\frac{m^2-2}{d}\right)\in\Sq.
\end{equation}
If the square in \eqref{eq:conditions} is zero, the two choices of $\delta$
coincide; otherwise the parametrization is unique after fixing the sign of
$z$.
\end{theorem}

\begin{proof}
If $x=0$, the equation becomes $y^2=2$, impossible in $\mathbb Z$.  Now write
$x=\varepsilon d$, where $d=|x|>0$ and $\varepsilon=\operatorname{sgn}(x)$,
and put $m=y$.  Since $x^2=d^2$, the equation is
\[
 m^2+d^2m+\varepsilon dz^2=2.
\]
Reduction modulo $d$ gives $d\mid m^2-2$, and rearrangement gives
\[
 z^2=-\varepsilon\left(dm+\frac{m^2-2}{d}\right).
\]
Thus every solution has the form \eqref{eq:classification}--\eqref{eq:conditions}.
Conversely, substituting those expressions yields
\[
 m^2+d^2m+\varepsilon d
 \left[-\varepsilon\left(dm+\frac{m^2-2}{d}\right)\right]
 =m^2-(m^2-2)=2.
\]
Hence every listed triple is a solution, and the recovery
$d=|x|$, $\varepsilon=\operatorname{sgn}(x)$, $m=y$ proves completeness and
the stated uniqueness.
\end{proof}

The classification is unconditional and effective: for any prescribed $y=m$,
one checks only the finitely many positive divisors of $m^2-2$ and tests a
single integer for being a square.

\section{An explicit infinite family}

The quadratic in $y$ is best centred by
\begin{equation}\label{eq:normidentity}
 (2y+x^2)^2+4xz^2=x^4+8,
\end{equation}
which follows by multiplying the original equation by $4$ and completing the
square.

\begin{proposition}[A Pell family on $x=-17$]\label{prop:family}
Define integer sequences $(U_n)$ and $(Z_n)$ by
\begin{equation}\label{eq:recurrence}
 U_0=573,\quad Z_0=60,\qquad
 \begin{cases}
 U_{n+1}=33U_n+272Z_n,\\
 Z_{n+1}=4U_n+33Z_n
 \end{cases}
 \quad(n\ge0).
\end{equation}
Then, for every $n\ge0$ and every $\sigma,\tau\in\{\pm1\}$,
\begin{equation}\label{eq:family}
 \boxed{
 (x,y,z)=
 \left(-17,\frac{\sigma U_n-289}{2},\tau Z_n\right)
 }
\end{equation}
is an integral solution.  These give infinitely many distinct solutions.
\end{proposition}

\begin{proof}
In $\mathbb Z[\sqrt{17}]$, recurrence \eqref{eq:recurrence} is exactly
\[
 U_{n+1}+2Z_{n+1}\sqrt{17}
 =(33+8\sqrt{17})(U_n+2Z_n\sqrt{17}).
\]
Since $33^2-17\cdot8^2=1$, norms are preserved.  The initial norm is
\[
 573^2-68\cdot60^2=83529=17^4+8,
\]
so
\begin{equation}\label{eq:invariant}
 U_n^2-68Z_n^2=17^4+8
\end{equation}
for every $n$.  Moreover $U_0$ is odd and \eqref{eq:recurrence} preserves
oddness of $U_n$, so the displayed $y$ is integral.  For $x=-17$,
\eqref{eq:normidentity} is
\[
 (2y+289)^2-68z^2=17^4+8;
\]
substitution of \eqref{eq:family} and \eqref{eq:invariant} proves the claim.
Finally, $U_n,Z_n>0$ and $U_{n+1}>U_n$ by induction, so the resulting
solutions are pairwise distinct as $n$ varies.
\end{proof}

The first two Vieta pairs are
\[
 (-17,142,\pm60),\quad (-17,-431,\pm60),
\]
\[
 (-17,17470,\pm4272),\quad (-17,-17759,\pm4272).
\]

\section{How the construction was found}

\textbf{1. Use the equation's quadratic structure.}
For fixed $x,z$, the equation is the monic quadratic
\[
 y^2+x^2y+(xz^2-2)=0.
\]
Hence the second root corresponding to a solution $y$ is $-x^2-y$.  This
Vieta symmetry suggests translating the variable to the centre
$-x^2/2$, namely setting $U=2y+x^2$.  One then obtains
\[
 U^2+4xz^2=x^4+8.
\]
The same equation also explains the complete classification: reducing the
original equation modulo $x$ gives $x\mid y^2-2$, after which solving
linearly for $z^2$ gives Theorem~\ref{thm:classification} in one step.

\textbf{2. The sign of $x$ reveals where infinitude can occur.}
If $x>0$, then $U^2+4xz^2=x^4+8$ is positive definite; for a fixed $x$ it
admits only finitely many pairs $(U,z)$.  Put instead $x=-a<0$.  The equation
becomes
\begin{equation}\label{eq:pellform}
 U^2-a(2z)^2=a^4+8,
\end{equation}
which is a generalized Pell equation.  Thus one seed solution can be
multiplied repeatedly by a norm-one unit to produce infinitely many others.
The task is therefore to choose $a$ so that both the unit and the seed are
simple.

\textbf{3. Engineer a simple unit and a compatible seed.}
A particularly transparent norm-one unit appears when
$a=r^2+1$, because
\[
 (2a-1)^2-a(2r)^2=1.
\]
To seek a low-complexity seed for \eqref{eq:pellform}, take
$U=2a^2-(r+1)$.  This automatically gives an integral
$y=(U-a^2)/2$ when $a$ and $r+1$ have the same parity.  Modulo $a$, the
required divisibility of
\[
 U^2-(a^4+8)=3a^4-4(r+1)a^2+(r+1)^2-8
\]
reduces, using $a=r^2+1$, to the very small condition
\[
 (r+1)^2-8\equiv 2r-8\pmod a.
\]
The clean choice is therefore $r=4$, which gives
\[
 a=17,\qquad U=2\cdot17^2-5=573.
\]
The remaining coordinate is then forced:
\[
 z^2=\frac{573^2-(17^4+8)}{4\cdot17}=3600,
 \qquad z=60,
\]
and consequently $y=(573-289)/2=142$.  Thus the seed is not a mysterious
computer output: it is produced by matching a simple Pell unit with a
small-offset ansatz for $U$.

\textbf{4. Propagate the seed.}
For $r=4$, the unit above is
$33+8\sqrt{17}$.  Multiplying
$573+120\sqrt{17}$ by successive powers of this unit preserves its norm
$17^4+8$ and gives exactly recurrence \eqref{eq:recurrence}.  Changing the
sign of the square-root coefficient changes $z$ to $-z$, while replacing
$U$ by $-U$ gives the Vieta partner $y\mapsto-17^2-y$.  These two elementary
symmetries account for all four sign choices in \eqref{eq:family}.

\end{document}
